\chapter*{Glossário}\addcontentsline{toc}{chapter}{GLOSSÁRIO}\label{ap:glossario}
\begin{description}[style=nextline,leftmargin=0pt,labelindent=0pt]
\item[Ablação.] Alteração ou remoção de um componente para comparar seu efeito dentro de um desenho definido.
\item[Agente.] Programa ao qual uma tarefa pode ser atribuída; o estudo com cartões calcula rankings e não executa os agentes.
\item[Arrependimento.] Diferença entre a melhor utilidade do oráculo e a utilidade da escolha realizada.
\item[Axioma.] Proposição admitida pelo sistema formal. A auditoria registra \texttt{propext}, \texttt{Classical.choice} e \texttt{Quot.sound}, sem demonstrar minimalidade desse conjunto.
\item[Calibração probabilística.] Correspondência entre probabilidades previstas e frequências observadas; não é garantida apenas pela soma unitária do softmax.
\item[Combinação convexa.] Soma ponderada com pesos não negativos de soma unitária.
\item[Comparador.] Estratégia utilizada como referência para avaliação no mesmo desenho.
\item[Intervalo exploratório.] Faixa calculada pela aproximação normal registrada, sem correção para agrupamento nas piscinas deste estudo.
\item[Maximizador.] Elemento de conjunto finito não vazio que atinge o maior valor de uma função; T3 exige finitude e não-vazio como hipóteses.
\item[Oráculo.] Critério sintético usado para avaliar as escolhas; não representa qualidade observada em tarefas reais.
\item[Overflow.] Excesso da faixa representável em ponto flutuante. T11 prova um limite da exponencial em reais exatos e não certifica todas as operações da implementação.
\item[Reprodução.] Execução dos procedimentos por outro executor a partir dos artefatos e do ambiente especificado; reprodução independente permanece pendente.
\item[Roteador.] Mecanismo que filtra candidatos e ordena suas utilidades para selecionar um agente.
\item[Semente.] Valor que inicializa uma sequência pseudoaleatória.
\item[Softmax.] Normalização exponencial de uma família finita não vazia de escores. Em reais exatos, seus pesos somam um e são positivos; em ponto flutuante, underflow pode produzir zeros.
\item[Temperatura.] Parâmetro positivo que controla a concentração dos pesos. T10 utiliza pelo menos dois candidatos e maximizador único como hipóteses suficientes.
\item[Top-\textit{k}.] Prefixo de tamanho $k$ de um ranking; T9 pressupõe a ordenação decrescente da lista.
\item[Underflow.] Resultado muito pequeno para ser representado com a precisão ou a faixa disponíveis.
\item[Verificação por máquina.] Checagem de um termo de prova pelo núcleo do assistente, relativa ao enunciado e aos axiomas utilizados; não garante que o enunciado modele integralmente o sistema concreto.
\end{description}
